\section{Related Work}
\label{sec:related}

\subsection{Weight-Space Model Merging}
Weight merging has emerged as a dominant paradigm for combining multiple specialized models fine-tuned from a common pre-trained base, avoiding the prohibitive cost of multi-task training or joint fine-tuning \cite{wortsman2022model, ilharco2022editing}. Task Arithmetic \cite{ilharco2022editing} demonstrates that task-specific fine-tuning can be represented as a task vector (the difference between the fine-tuned and pre-trained weights) and that adding these task vectors to the base pre-trained model yields an effective multi-task model.
Subsequent research has identified and addressed several limitations of naive linear merging. For instance, TIES-Merging \cite{yadav2023ties} addresses sign disagreements and redundant parameters in task vectors by resolving conflicts and sparsifying weights. DARE \cite{yu2023language} uses random dropping and scaling to merge models, showing that a high proportion of weights can be zeroed out without losing task expert capabilities. 
While these methods focus on pruning and sign correction, they still heavily rely on manual or heuristically determined merging coefficients, often choosing uniform weights across all layers.

\subsection{Test-Time Adaptation (TTA) in Model Merging}
To overcome the limitations of static merging coefficients, several authors have proposed adjusting coefficients dynamically at test-time. This line of research is inspired by classical Test-Time Adaptation (TTA) in computer vision, which adapts a model to target distribution shifts on unlabeled test streams using unsupervised objectives like entropy minimization \cite{wang2020tent, liang2023source}.
AdaMerging \cite{adamerging} was the first to apply this to model merging. It proposes both task-wise and layer-wise variants, optimizing coefficients on-the-fly using Adam on the unlabeled test stream to minimize the entropy of the merged model's predictions. RegCalMerge \cite{regcalmerge} extends this by introducing class-capacity normalization and elastic spatial regularization to prevent representational collapse. PolyMerge \cite{polymerge} constraints the search space of merging coefficients by modeling them as a continuous polynomial function of depth, optimizing the polynomial parameters on-the-fly. Q-Merge \cite{qmerge} adapts model merging to quantized networks using straight-through estimators.
All of these methods report substantial improvements over the naive uniform baseline, claiming that online test-time optimization is essential for resolving representational interference and adapting to local stream statistics.

\subsection{Methodological Skepticism and Baselines}
The field of machine learning has a long history of "illusionary progress" driven by weak or unoptimized baselines, flawed evaluation protocols, and a lack of rigorous robustness testing \cite{bulo2017loss, musgrave2020metric, liao2021are}. For example, in metric learning, Musgrave et al.\ \cite{musgrave2020metric} demonstrated that standard, well-tuned classical baselines often outperform highly complex modern architectures when evaluated under identical, fair conditions. Similarly, in robustness and domain generalization, Gulrajani \& Lopez-Paz \cite{gulrajani2020in} showed that simple empirical risk minimization is extremely hard to beat when hyperparameter tuning is conducted rigorously.
Our work adopts this skeptical, methodologically rigorous stance towards the online TTA model merging literature. We hypothesize that the reported benefits of TTA are overemphasized due to the choice of a weak, unoptimized baseline (uniform merging) as the sole competitor (the "no-data" strawman). Furthermore, we argue that evaluating TTA methods solely on stable, clean, i.e.\ i.i.d.\ test streams obscures their fundamental fragility. By introducing OFS-Tune---a simple, offline, few-shot tuned baseline---and stress-testing all methods under adversarial stream shifts, we aim to provide the community with a much-needed methodological reality check.
